\section{Pattern 4：Inversion（反转）}

Agent 天生倾向于\textbf{猜测并立即生成}。Inversion 模式翻转了这一动态：不是用户驱动 prompt、Agent 执行，而是 \textbf{Agent 充当采访者}，在行动之前先收集完整的上下文。

\begin{importantbox}{Inversion 的核心约束}
Inversion 依赖于明确的、不可协商的门控指令（gating instructions），例如"在所有阶段完成之前，不得开始构建"。Agent 按顺序提出结构化问题，等待用户回答后才进入下一阶段，直到获得完整的需求和部署约束后才生成最终输出。
\end{importantbox}

\subsection{工作原理}

\begin{enumerate}
    \item Agent 按预定义的阶段依次向用户提问
    \item 每个阶段收集特定类别的信息（需求、约束、偏好等）
    \item 严格的门控条件阻止 Agent 跳过任何阶段
    \item 只有所有答案收集完毕后，Agent 才综合生成最终输出
\end{enumerate}

\begin{warningbox}{为什么需要显式门控？}
如果没有显式的门控指令，Agent 很容易在收集到部分信息后就急于生成结果。这会导致输出质量低下，遗漏关键需求。门控指令（如 ``DO NOT start building until all phases are complete''）是 Inversion 模式的生命线，必须用强硬、明确的语言表达。
\end{warningbox}

\subsection{适用场景}

\begin{itemize}
    \item 项目规划（需要全面了解需求后才能制定方案）
    \item 架构设计（需要了解各种约束条件）
    \item 复杂任务分解（需要明确上下文后才能合理拆分）
\end{itemize}

\subsection{本章小结}

Inversion 模式通过显式门控指令，强制 Agent 在行动前完成信息收集。这种"先问后做"的策略显著提升了输出质量，特别适合需要全面上下文的复杂任务。

